% ============================================================
%  Relatório — Pipeline AIH/SUS + Dashboard Streamlit
% ============================================================
\documentclass[12pt,a4paper]{report}

\usepackage[utf8]{inputenc}
\usepackage[T1]{fontenc}
\usepackage[brazil]{babel}
\usepackage{geometry}
\geometry{a4paper, margin=2.5cm}

\usepackage{hyperref}
\usepackage{listings}
\usepackage{xcolor}
\usepackage{enumitem}
\usepackage{booktabs}

% ── Configuração do listings ──────────────────────────────────────────────────
\definecolor{codebg}{rgb}{0.95,0.95,0.95}
\definecolor{codeframe}{rgb}{0.7,0.7,0.7}
\lstset{
  basicstyle=\small\ttfamily,
  backgroundcolor=\color{codebg},
  frame=single,
  rulecolor=\color{codeframe},
  breaklines=true,
  captionpos=b,
  numbers=left,
  numberstyle=\tiny\color{gray},
  stepnumber=1,
  tabsize=4,
  showstringspaces=false,
}
\lstdefinestyle{sql}{language=SQL, morekeywords={RECURSIVE,WITH,UNION,ALL}}

% ─────────────────────────────────────────────────────────────────────────────
\begin{document}

\title{Relatório Técnico: Pipeline de Dados AIH/SUS e\\
       Dashboard Interativo com Streamlit}
\author{Caio Resende Junqueira}
\date{2025}
\maketitle
\tableofcontents
\newpage

% =============================================================
\chapter{Desenvolvimento da Aplicação Web com Streamlit}
% =============================================================

Com os dados devidamente migrados e estruturados no servidor PostgreSQL do IESB, o
passo subsequente consistiu na construção de uma camada de apresentação e análise
interativa: um \textit{dashboard} web desenvolvido com o \textit{framework}
\textbf{Streamlit}, hospedado na plataforma \textbf{Streamlit Cloud} e integrado ao
banco de dados via protocolo TCP/IP.

% -------------------------------------------------------------
\section{Criação do Repositório e Estrutura do Projeto}
% -------------------------------------------------------------

O código-fonte da aplicação foi versionado e publicado numa plataforma de controlo de
versões distribuída (\textit{Git/GitHub}), sob o repositório
\texttt{caiorj2004/aih\_sp}. A adoção do controlo de versões é uma prática
arquitetural indispensável, pois permite a rastreabilidade de cada alteração, a
reversão a estados anteriores estáveis e, fundamentalmente, a integração direta com
serviços de hospedagem como o Streamlit Cloud, que realiza o \textit{deploy}
automaticamente a partir de um \textit{branch} designado do repositório.

A estrutura de diretórios foi definida de forma minimalista, separando
responsabilidades em ficheiros distintos:

\begin{verbatim}
aih_sp/
├── app.py                # Interface Streamlit (camada de apresentação)
├── db.py                 # Conexão e consultas ao PostgreSQL (camada de dados)
├── fallback.py           # Acesso a dados locais Parquet (modo offline)
├── requirements.txt      # Dependências Python declarativas
├── data/
│   ├── aih_qtd_fallback.parquet  # Snapshot de quantidades
│   └── aih_vl_fallback.parquet   # Snapshot de valores
├── .streamlit/
│   ├── config.toml       # Tematização e configurações de servidor (versionado)
│   └── secrets.toml      # Credenciais de acesso — NÃO versionado
└── .gitignore
\end{verbatim}

A organização adota três módulos funcionais com responsabilidades claras, seguindo o
princípio de responsabilidade única (\textit{Single Responsibility Principle}):

\begin{itemize}
    \item \texttt{app.py} — toda a lógica de interface e renderização;
    \item \texttt{db.py} — toda a comunicação com o SGBD: consultas SQL, gestão de
    conexões e \textit{cache};
    \item \texttt{fallback.py} — leitura dos arquivos Parquet locais utilizados
    quando o banco de dados não está acessível.
\end{itemize}

Esta separação facilita a manutenção independente das camadas e permite que eventuais
substituições de \textit{backend} exijam alterações apenas no módulo de dados, sem
tocar na camada de apresentação.

\subsection{Proteção de Credenciais e \texttt{.gitignore}}

Um risco crítico no desenvolvimento de aplicações que consomem serviços remotos é a
exposição acidental de credenciais em repositórios públicos. Para mitigar este risco, o
ficheiro \texttt{.streamlit/secrets.toml} — que contém \textit{host}, porta, nome do
banco, utilizador e senha do PostgreSQL — foi explicitamente declarado no
\texttt{.gitignore}, garantindo que nenhuma credencial sensível seja comprometida
em nenhum \textit{commit}.

Para o ambiente de produção (Streamlit Cloud), as credenciais são injetadas de forma
segura através do painel da plataforma em \textbf{Settings → Secrets}, sem nunca
transitarem pelo sistema de controlo de versões.

% -------------------------------------------------------------
\section{Declaração de Dependências}
% -------------------------------------------------------------

As dependências da aplicação foram declaradas no ficheiro \texttt{requirements.txt},
que é processado automaticamente pelo Streamlit Cloud durante o provisionamento do
ambiente de execução. As bibliotecas e as restrições de versão mínima adotadas foram:

\begin{itemize}
    \item \textbf{\texttt{streamlit>=1.33}} — \textit{framework} de interface web,
    responsável pela renderização de componentes reativos, gestão de estado de sessão e
    mecanismo nativo de conexões de dados.
    \item \textbf{\texttt{pandas>=2.0}} — biblioteca de manipulação e análise de dados
    tabulares; utilizada para transformações pós-consulta e merge entre DataFrames.
    \item \textbf{\texttt{plotly>=5.20}} — biblioteca de visualização interativa;
    responsável pelos gráficos de linha de duplo eixo, barras de ranking,
    dispersão, \textit{treemap} e \textit{heatmap}.
    \item \textbf{\texttt{sqlalchemy>=2.0}} — \textit{toolkit} de abstração de banco de
    dados; fornece a interface de dialeto SQL e o mecanismo de execução de consultas
    parametrizadas, atuando como adaptador entre a aplicação Python e o driver
    PostgreSQL.
    \item \textbf{\texttt{psycopg2-binary>=2.9}} — \textit{driver} DBAPI nativo do
    PostgreSQL para Python; responsável pela comunicação de baixo nível via socket TCP.
    A variante \textit{binary} inclui os binários compilados estaticamente,
    eliminando a dependência de bibliotecas do sistema operacional no servidor de
    hospedagem.
    \item \textbf{\texttt{pyarrow}} — motor de leitura de ficheiros Parquet utilizado
    pelo Pandas no modo \textit{fallback}.
\end{itemize}

% -------------------------------------------------------------
\section{Módulo de Dados: \texttt{db.py}}
% -------------------------------------------------------------

O módulo \texttt{db.py} constitui a espinha dorsal da integração entre a aplicação e o
banco de dados relacional. A sua arquitetura abrange quatro preocupações principais:
gestão de conexão, segurança das consultas, introspecção dinâmica de esquema e
otimização de desempenho.

\subsection{Conexão ao PostgreSQL via Streamlit Secrets}

A conexão ao servidor PostgreSQL é estabelecida através do mecanismo nativo de
conexões do Streamlit (\texttt{st.connection}), que lê automaticamente o bloco
\texttt{[connections.postgresql]} do ficheiro de \textit{secrets}. O objeto de
conexão retornado encapsula internamente um \textit{engine} SQLAlchemy com
\textit{pool} de conexões, tornando o processo transparente para a camada de consulta.

Para evitar que múltiplas execuções do script Streamlit — provocadas pela natureza
reativa do \textit{framework}, que re-executa o script a cada interação do utilizador
— criem conexões redundantes ao banco de dados, a função de conexão foi decorada com
\texttt{@st.cache\_resource}. Este decorator instrui o Streamlit a manter o objeto de
conexão como um \textit{singleton} partilhado pela sessão e por todos os utilizadores
simultâneos:

\begin{lstlisting}[language=Python]
@st.cache_resource
def get_connection():
    return st.connection("postgresql", type="sql")
\end{lstlisting}

\subsection{Estratégias de Cache e Invalidação por TTL}

Num ambiente de \textit{dashboard} analítico, as mesmas consultas são disparadas
repetidamente a cada interação do utilizador. Para evitar sobrecarga desnecessária no
SGBD, todas as funções de consulta foram decoradas com \texttt{@st.cache\_data}, que
armazena o resultado em memória e invalida o cache após um \textit{Time-To-Live} (TTL)
configurável:

\begin{itemize}
    \item \textbf{TTL de 3600 segundos (1 hora):} aplicado às funções que consultam
    metadados de esquema (\texttt{get\_table\_columns} e
    \texttt{get\_dimension\_mapping}), cujo conteúdo raramente varia em produção.
    \item \textbf{TTL de 1800 segundos (30 minutos):} aplicado às funções de
    carregamento de opções de filtros (\texttt{load\_filter\_options} e
    \texttt{load\_municipality\_options}) e de dados consolidados
    (\texttt{load\_consolidated\_data}), refletindo a frequência típica de atualização
    dos dados hospitalares.
    \item \textbf{TTL de 900 segundos (15 minutos):} aplicado às consultas de totais
    por período (\texttt{get\_period\_totals}) e às colunas de procedimentos
    (\texttt{get\_procedure\_columns}), onde maior frescor dos dados é desejável.
\end{itemize}

\subsection{Segurança: \textit{Whitelist} de Tabelas e Consultas Parametrizadas}

A aplicação exige proteção contra ataques de injeção de SQL (\textit{SQL Injection}),
visto que nomes de tabelas e colunas são construídos dinamicamente a partir de dados
retornados pelo próprio banco de dados. Para isso, adotaram-se duas defesas em camadas:

\subsubsection*{Whitelist de Tabelas Permitidas}
Uma constante \texttt{\_ALLOWED\_TABLES} define o conjunto exato de tabelas que podem
ser consultadas. Qualquer chamada à função \texttt{get\_table\_columns()} com um nome
fora deste conjunto resulta em \texttt{ValueError} imediato, antes de qualquer acesso
ao banco:

\begin{lstlisting}[language=Python]
_ALLOWED_TABLES = {"aih_qtd", "aih_vl", "municipios_ibge", "unidade_federacao"}

def get_table_columns(table_name: str) -> List[str]:
    if table_name not in _ALLOWED_TABLES:
        raise ValueError(f"Tabela '{table_name}' nao esta na lista de tabelas permitidas.")
    ...
\end{lstlisting}

\subsubsection*{Consultas Parametrizadas com Parâmetros Nomeados}
Todos os valores dinâmicos fornecidos pelo utilizador (anos, meses, códigos de UF,
códigos de município) são injetados nas consultas SQL exclusivamente através de
parâmetros nomeados do SQLAlchemy (notation \texttt{:param\_name}), nunca por
concatenação de strings. A função utilitária \texttt{\_build\_in\_clause()} gera
automaticamente as listas de placeholders e popula o dicionário de parâmetros:

\begin{lstlisting}[language=Python]
def _build_in_clause(column_sql, values, prefix, params):
    placeholders = []
    for index, value in enumerate(values):
        key = f"{prefix}_{index}"
        params[key] = value
        placeholders.append(f":{key}")
    return f" AND {column_sql} IN ({', '.join(placeholders)}) "
\end{lstlisting}

Esta abordagem garante que o \textit{driver} DBAPI trate todos os valores como literais
escapados, tornando impossível a modificação da estrutura sintática da instrução SQL
por parte de um atacante.

\subsection{Introspecção Dinâmica de Esquema}

A função \texttt{get\_dimension\_mapping()} implementa um padrão de descoberta dinâmica
das colunas das tabelas de dimensão, consultando o \texttt{information\_schema.columns}
do PostgreSQL. Este mecanismo torna a aplicação resiliente a variações na nomenclatura
das colunas entre diferentes versões da carga de dados — por exemplo, a coluna com o
nome do município pode estar nomeada como \texttt{nome\_municipio}, \texttt{no\_municipio}
ou simplesmente \texttt{municipio}. A função \texttt{\_first\_existing()} itera sobre
uma lista de candidatos em ordem de preferência e retorna o primeiro que existir no
esquema real:

\begin{lstlisting}[language=Python]
municipio_name_col = _first_existing(
    ["nome_municipio", "no_municipio", "municipio", "nome"],
    municipio_cols,
    municipio_code_col,   # fallback
)
\end{lstlisting}

\subsection{Otimização de Desempenho: CTE Recursiva \textit{Server-Side}}

A função \texttt{load\_filter\_options()} é responsável por popular os seletores de
filtro do \textit{dashboard} na inicialização da aplicação. A implementação ingênua
deste requisito utilizaria \texttt{SELECT DISTINCT ano FROM aih\_qtd}, que força o
PostgreSQL a varrer a totalidade da tabela fato — uma operação de complexidade O(n),
onde n é a cardinalidade da tabela.

Para contornar este gargalo, foi adotado o padrão de \textit{``Loose Index Scan''}
implementado via \textbf{CTE Recursiva} (\textit{Common Table Expression}). Esta
técnica aproveita a estrutura de B-tree do índice da coluna para saltar diretamente de
um valor distinto ao seguinte, executando apenas k iterações (onde k é o número de
valores distintos, geralmente 5 a 12 para o campo ano e exatamente 12 para o campo
mês), com custo logarítmico em n por iteração:

\begin{lstlisting}[style=sql]
WITH RECURSIVE t(ano) AS (
    (SELECT ano FROM aih_qtd WHERE ano IS NOT NULL ORDER BY ano LIMIT 1)
    UNION ALL
    SELECT (
        SELECT ano FROM aih_qtd
        WHERE ano > t.ano AND ano IS NOT NULL
        ORDER BY ano LIMIT 1
    )
    FROM t WHERE t.ano IS NOT NULL
)
SELECT ano FROM t WHERE ano IS NOT NULL
\end{lstlisting}

Adicionalmente, a execução destas consultas foi realizada através do contexto
\texttt{conn.session} (interface de sessão bruta do SQLAlchemy) com as opções de
execução \texttt{stream\_results=True} e \texttt{yield\_per=50}. Este mecanismo
instrui o PostgreSQL a manter o cursor aberto no servidor e a transmitir os resultados
em lotes de 50 linhas (\textit{server-side cursor}), evitando a materialização da
totalidade do resultado em memória antes do processamento:

\begin{lstlisting}[language=Python]
with conn.session as session:
    result = session.execute(
        _years_sql,
        execution_options={"stream_results": True, "yield_per": 50},
    )
    for batch in result.partitions(50):
        for row in batch:
            years_set.add(str(row[0]).strip())
\end{lstlisting}

\subsection{Consulta Consolidada em Dois Estágios}

A função \texttt{load\_consolidated\_data()} executa o JOIN principal entre as quatro
tabelas do modelo (\texttt{aih\_qtd}, \texttt{aih\_vl}, \texttt{municipios\_ibge} e
\texttt{unidade\_federacao}). Para evitar que o PostgreSQL materialize o produto
cartesiano das tabelas fato antes da aplicação dos filtros geográficos — operação que
esgotaria o espaço temporário em disco para conjuntos de dados de grande cardinalidade
— a consulta foi decomposta em dois estágios sequenciais:

\begin{enumerate}
    \item \textbf{Resolução das dimensões geográficas:} uma consulta prévia,
    exclusivamente sobre as tabelas de dimensão (\texttt{municipios\_ibge} e
    \texttt{unidade\_federacao}), retorna a lista de \texttt{cod\_municipio} válidos
    para as UFs e municípios selecionados pelo utilizador.
    \item \textbf{Consulta às tabelas fato:} a lista de códigos obtida na etapa anterior
    é utilizada como filtro direto na cláusula \texttt{WHERE} da consulta sobre
    \texttt{aih\_qtd} e \texttt{aih\_vl}, permitindo que o PostgreSQL utilize os índices
    sobre \texttt{cod\_municipio} e evite \textit{full table scans}.
\end{enumerate}

A junção final entre os dados de facto e as informações de dimensão é realizada em
Python via \texttt{DataFrame.merge()}, após a materialização dos dois conjuntos de
resultados, garantindo a eficiência da comunicação com o SGBD.

% -------------------------------------------------------------
\section{Módulo de Fallback: \texttt{fallback.py}}
% -------------------------------------------------------------

Durante o desenvolvimento verificou-se que a dependência exclusiva do servidor
PostgreSQL tornava a aplicação vulnerável a indisponibilidades de rede ou de
infraestrutura. Para garantir resiliência operacional, foi desenvolvido o módulo
\texttt{fallback.py}, que implementa um modo \textit{offline} completo baseado em
ficheiros \textbf{Parquet} armazenados localmente na pasta \texttt{data/} do
repositório.

\subsection{Ficheiros Parquet}

Dois ficheiros Parquet servem de fonte de dados alternativa:

\begin{itemize}
    \item \texttt{data/aih\_qtd\_fallback.parquet} — colunas de quantidades por
    procedimento: \texttt{ano}, \texttt{mes}, \texttt{cod\_municipio},
    \texttt{municipio}, \texttt{qtd\_*}, \texttt{total}.
    \item \texttt{data/aih\_vl\_fallback.parquet} — colunas de valores por
    procedimento: \texttt{ano}, \texttt{mes}, \texttt{cod\_municipio},
    \texttt{municipio}, \texttt{vl\_*}, \texttt{total}.
\end{itemize}

Uma limitação importante: os ficheiros Parquet foram gerados a partir das exportações
brutas do TabNet, que não inclui a dimensão de Unidade da Federação (UF). Por essa
razão, o modo \textit{offline} opera apenas com granularidade municipal.

\subsection{Funções do Módulo}

O módulo \texttt{fallback.py} oferece um conjunto de funções espelhadas às do
\texttt{db.py}, com assinatura compatível para permitir substituição transparente:

\begin{itemize}
    \item \texttt{load\_fallback\_raw()} — lê e cacheia os dois ficheiros Parquet com
    \texttt{@st.cache\_data}, normalizando os tipos das colunas-chave (\texttt{ano},
    \texttt{mes}, \texttt{cod\_municipio}) para \texttt{str}.
    \item \texttt{get\_fallback\_filter\_options()} — extrai anos, meses e o DataFrame
    de municípios disponíveis no Parquet.
    \item \texttt{load\_fallback\_consolidated()} — filtra e mescla os dois Parquets
    em função dos filtros de período e município, retornando a mesma estrutura de
    DataFrame que \texttt{load\_consolidated\_data()}.
    \item \texttt{get\_fallback\_period\_totals()} — equivalente offline de
    \texttt{get\_period\_totals()}, usado no cálculo dos deltas dos KPIs.
    \item \texttt{get\_fallback\_procedure\_columns()} — retorna as listas de colunas
    \texttt{qtd\_*} e \texttt{vl\_*} disponíveis nos Parquets.
\end{itemize}

\subsection{Ativação Automática e Modo Misto}

O processo de ativação do modo \textit{offline} ocorre em duas vias:

\begin{enumerate}
    \item \textbf{Ativação automática por falha:} se a tentativa de conexão ao
    PostgreSQL lançar qualquer exceção (\texttt{OperationalError}, timeout ou erro de
    autenticação), o módulo \texttt{app.py} captura o erro e, imediatamente,
    tenta carregar os ficheiros Parquet. Se bem-sucedido, o \textit{flag}
    \texttt{\_using\_fallback} é definido como \texttt{True} e a aplicação continua
    operando em modo \textit{offline}:

\begin{lstlisting}[language=Python]
try:
    _years, _months, _uf_options = load_filter_options()
except Exception as exc:
    _db_error = exc
    try:
        _fb_years, _fb_months, _fb_municipios = get_fallback_filter_options()
        _years, _months = _fb_years, _fb_months
        _using_fallback = True
    except Exception:
        pass  # both DB and fallback unavailable
\end{lstlisting}

    \item \textbf{Toggle manual na barra lateral:} mesmo quando o banco de dados está
    disponível, um botão de alternância (\texttt{st.toggle}) é exibido na barra lateral
    — desde que os ficheiros Parquet também estejam presentes. Este controlo permite
    ao utilizador alternar voluntariamente entre os dados do servidor e os dados locais,
    o que é útil para comparar versões dos dados ou para demonstrações
    sem acesso à rede:

\begin{lstlisting}[language=Python]
if _db_error is None and _fb_municipios is not None:
    if st.toggle("Usar dados locais (Parquet)", key="use_local"):
        _using_fallback = True
        _years = _fb_years
        _months = _fb_months
\end{lstlisting}
\end{enumerate}

Se ambas as fontes falharem, a função auxiliar \texttt{\_render\_db\_unavailable()}
exibe uma mensagem de erro detalhada, o \textit{traceback} completo da exceção do banco
e interrompe a renderização das abas de dados, evitando que o utilizador receba apenas
um ecrã em branco.

\subsection{Adaptações da Interface no Modo Offline}

Quando o \textit{flag} \texttt{\_using\_fallback} está ativo, diversas componentes da
interface são automaticamente adaptadas para refletir a ausência da dimensão UF:

\begin{itemize}
    \item O filtro \textbf{Unidade da Federação} é removido da barra lateral; apenas
    o filtro de Município permanece disponível.
    \item Uma faixa de aviso é exibida:~⚠️~\textbf{Modo Offline} (se o banco falhou)
    ou~📁~informativo (se o utilizador activou o toggle manualmente).
    \item No gráfico \textbf{Ranking Top 10}, o rádio de seleção de nível
    (Município / UF) é ocultado; o ranking opera exclusivamente por Município.
    \item No \textbf{Scatter Plot}, o \textit{hover} exibe apenas o nome do município,
    sem a coluna de UF.
    \item No \textbf{Heatmap Sazonal}, a opção de eixo Y por UF é suprimida.
    \item Na aba de \textbf{Estatísticas Descritivas}, o indicador
    ``UFs no recorte'' é omitido.
\end{itemize}

Os demais recursos — série temporal, treemap por categoria, KPIs com delta e
exportação CSV — funcionam sem alterações.

% -------------------------------------------------------------
\section{Módulo de Interface: \texttt{app.py}}
% -------------------------------------------------------------

O módulo \texttt{app.py} implementa a totalidade da camada de apresentação do
\textit{dashboard}, incluindo a configuração da página, a gestão de filtros
hierárquicos, o cálculo e exibição de KPIs, as visualizações analíticas interativas
e o dicionário de variáveis.

\subsection{Configuração e Tematização}

A identidade visual da aplicação foi configurada no ficheiro
\texttt{.streamlit/config.toml}, definindo uma paleta de cores corporativa (azul
institucional \texttt{\#0066CC} como cor primária, fundo branco e texto escuro
\texttt{\#1A1A2E}), além de desativar a telemetria de uso do Streamlit:

\begin{lstlisting}[numbers=none]
[theme]
primaryColor             = "#0066CC"
backgroundColor          = "#FFFFFF"
secondaryBackgroundColor = "#F0F2F6"
textColor                = "#1A1A2E"
font                     = "sans serif"

[server]
enableStatistics = false
\end{lstlisting}

A página foi configurada com layout largo (\texttt{layout="wide"}) para maximizar a
área de visualização dos gráficos, e o ícone e título da aba do navegador foram
personalizados via \texttt{st.set\_page\_config()}.

\subsection{Formatação Numérica Brasileira}

Para apresentar os valores monetários e contagens no padrão brasileiro (ponto como
separador de milhar, vírgula como separador decimal), foram implementadas duas funções
utilitárias:

\begin{itemize}
    \item \texttt{format\_currency(value)} — formata um float como moeda BRL:
    \texttt{R\$ 1.234,56}.
    \item \texttt{br\_format(value, decimals)} — formata qualquer número com o locale
    brasileiro; utilizada para renderizar o DataFrame de dados brutos e a tabela de
    estatísticas.
\end{itemize}

Ambas as funções utilizam o método \texttt{str.maketrans} para permutar pontos e
vírgulas após a formatação padrão do Python, evitando dependência de bibliotecas de
locale do sistema operacional.

\subsection{Dicionário de Rótulos (\texttt{\_PROC\_NAMES} e \texttt{COLUMN\_LABELS})}

A aplicação mantém um dicionário estático \texttt{\_PROC\_NAMES} com mais de 80
entradas, mapeando os códigos numéricos de grupos de procedimentos SIGTAP
(ex.: \texttt{"0406"}) para as respectivas denominações legíveis em português
(ex.: \texttt{"Cirurgia do aparelho circulatório"}). Este dicionário alimenta o
dicionário de exibição \texttt{COLUMN\_LABELS}, que associa nomes de colunas como
\texttt{qtd\_0406} ao rótulo completo \texttt{"Qtd – 0406 Cirurgia do aparelho
circulatório"}. A função \texttt{col\_label(col)} é então utilizada em todos os
componentes Streamlit (\texttt{format\_func}, títulos de eixos, \textit{tooltips})
para substituir os nomes técnicos das colunas por descrições compreensíveis.

\subsection{Filtros Hierárquicos na Barra Lateral}

A barra lateral (\texttt{st.sidebar}) implementa um sistema de filtragem hierárquica
em quatro níveis, onde cada nível restringe as opções do nível subsequente:

\begin{enumerate}
    \item \textbf{Ano} (\texttt{st.multiselect}): carregado da função
    \texttt{load\_filter\_options()}, populado com os anos presentes em
    \texttt{aih\_qtd}.
    \item \textbf{Mês} (\texttt{st.multiselect}): igualmente carregado de
    \texttt{load\_filter\_options()}, com opções ordenadas cronologicamente via a
    função utilitária \texttt{month\_to\_num()}, que suporta abreviações inglesas
    e portuguesas bem como valores numéricos.
    \item \textbf{Unidade da Federação} (\texttt{st.multiselect}): renderizado com
    labels legíveis no formato \textit{``SP — São Paulo''}, mapeando internamente
    para o código de UF utilizado nas consultas SQL.
    \item \textbf{Município} (\texttt{st.multiselect}): carregado dinamicamente pela
    função \texttt{load\_municipality\_options()}, cujas opções variam em função das
    UFs selecionadas no nível superior.
\end{enumerate}

No modo \textit{offline} (fallback), o filtro de UF é suprimido e o filtro de
Município é preenchido diretamente a partir do Parquet, sem dependência do PostgreSQL.

\subsection{KPIs com Cálculo de Delta Interperiódico}

A secção de KPIs do \textit{dashboard} apresenta três métricas de alto nível num
layout de três colunas (\texttt{st.columns(3)}):

\begin{itemize}
    \item \textbf{Total de Procedimentos:} soma de \texttt{total\_qtd} sobre todos os
    registos do recorte selecionado, formatado com separador de milhar brasileiro.
    \item \textbf{Valor Total Repassado:} soma de \texttt{total\_vl}, formatado no
    padrão monetário BRL via \texttt{format\_currency()}.
    \item \textbf{Valor Médio por Procedimento (\textit{Ticket} Médio):} quociente
    entre valor total e quantidade total, com tratamento seguro da divisão por zero via
    comparação com \texttt{EPSILON = 1e-9}.
\end{itemize}

Cada KPI exibe adicionalmente um \textbf{delta percentual} calculado em relação ao
período imediatamente anterior. A função \texttt{previous\_period()} implementa a
lógica de retrocesso calendárico, tratando corretamente as transições de ano
(dezembro → novembro do ano corrente, janeiro → dezembro do ano anterior). Os totais
do período anterior são consultados via \texttt{get\_period\_totals()} (ou
\texttt{get\_fallback\_period\_totals()} em modo offline), que aplica os mesmos filtros
geográficos da sessão atual para garantir comparabilidade.

% -------------------------------------------------------------
\section{Abas de Análise}
% -------------------------------------------------------------

O conteúdo do \textit{dashboard} é organizado em quatro abas
(\texttt{st.tabs}), separando categorias distintas de informação:

\subsection*{Aba 0 — 🏠 Introdução}

A aba de introdução é sempre renderizada, mesmo antes de qualquer filtro ser
selecionado, e não depende de conexão com o banco de dados. Apresenta três secções
principais:

\begin{itemize}
    \item \textbf{Origem dos Dados} — descreve o SIH/SUS, o portal TabNet e a
    estratégia de scraping com Playwright.
    \item \textbf{Pipeline de Engenharia de Dados} — resume as três etapas
    (Extração, Staging SQLite, Produção PostgreSQL) num layout de três colunas.
    \item \textbf{Modo Offline (Fallback)} — documenta ao utilizador final o
    comportamento da aplicação quando o banco de dados está indisponível, incluindo as
    funcionalidades que são adaptadas e as que permanecem intactas.
\end{itemize}

\subsection*{Aba A — Lista dos Dados Armazenados}

Apresenta o DataFrame consolidado resultante da consulta principal, com formatação
numérica brasileira aplicada coluna a coluna (0 casas decimais para contagens,
2 casas decimais para valores monetários). Um botão \texttt{st.download\_button}
permite a exportação imediata dos dados visíveis em formato CSV com codificação UTF-8.

Abaixo da tabela de dados é exibido um \textbf{Dicionário de Variáveis} dividido em
três tabelas: colunas de controlo e metadados, colunas de quantidades (\texttt{qtd\_*})
e colunas de valores financeiros (\texttt{vl\_*}), utilizando os rótulos do dicionário
\texttt{COLUMN\_LABELS}.

\subsection*{Aba B — Estatísticas Descritivas}

Aplica o método \texttt{DataFrame.describe().T} sobre as colunas métricas (totais e
colunas de procedimentos individuais) e enriquece o resultado com três estatísticas
adicionais: \textbf{variância} (\texttt{var()}), \textbf{assimetria}
(\texttt{skew()}) e \textbf{curtose} (\texttt{kurt()}), produzindo uma matriz
transposta completa formatada no locale brasileiro.

Três (ou dois, em modo \textit{offline}) indicadores de cardinalidade do recorte são
exibidos de forma destacada: número de municípios, número de UFs e número de períodos
(combinações únicas de ano/mês) no filtro ativo.

\subsection*{Aba C — Gráficos Analíticos}

Esta aba concentra cinco visualizações interativas construídas com a biblioteca
\textbf{Plotly}:

\begin{enumerate}
    \item \textbf{Série Temporal de Duplo Eixo:} gráfico de linhas com marcadores
    (\texttt{go.Scatter}) com dois eixos Y independentes e seletores de coluna para
    cada eixo. O eixo temporal é construído convertendo as colunas \texttt{ano} e
    \texttt{mes} em objetos \texttt{datetime} via \texttt{pd.to\_datetime()},
    ordenados cronologicamente, suportando os mesmos formatos de mês (numérico,
    abreviação inglesa ou abreviação portuguesa) tratados por \texttt{month\_to\_num()}.

    \item \textbf{Ranking Top 10:} gráfico de barras (\texttt{px.bar}) com seleção de
    nível de agregação (Município ou UF, este último indisponível em modo
    \textit{offline}) e de métrica. Sempre opera sobre todos os municípios da UF
    selecionada, independentemente do filtro de município ativo.

    \item \textbf{Scatter Plot de Procedimentos × Valor:} diagrama de dispersão
    (\texttt{px.scatter}) com seletores de eixo X e Y e escolha do coeficiente de
    correlação a exibir: \textbf{Pearson} (linear, sensível a outliers) ou
    \textbf{Spearman} (baseado em postos, robusto a outliers e não-paramétrico).
    A correlação de Spearman é calculada internamente via correlação de Pearson sobre
    os postos (\texttt{rank()}), eliminando a necessidade do pacote \texttt{scipy}.
    O valor do coeficiente é exibido abaixo do gráfico com uma nota metodológica.

    \item \textbf{Treemap por Categoria de Procedimento:} gráfico \textit{treemap}
    (\texttt{px.treemap}) substituindo o donut original. Permite selecionar as colunas
    de procedimento a incluir via \texttt{st.multiselect}; as colunas excluídas são
    agrupadas numa categoria ``Outros'' cujo valor é a soma das suas contribuições ao
    total do recorte. Os percentuais exibidos são sempre relativos ao total acumulado
    de todas as colunas disponíveis, garantindo consistência ao variar a seleção.

    \item \textbf{Heatmap Sazonal:} mapa de calor (\texttt{go.Heatmap}) com meses
    no eixo X e, no eixo Y, o ano ou a UF (configurável por rádio; UF indisponível em
    modo \textit{offline}). A escala de cores \texttt{YlOrRd} destaca visualmente
    os picos sazonais e quebras de padrão histórico. O eixo de meses usa abreviações
    em português (Jan–Dez) e opera sempre sobre todos os municípios da UF selecionada.
\end{enumerate}

% -------------------------------------------------------------
\section{Publicação no Streamlit Cloud}
% -------------------------------------------------------------

O \textit{deploy} da aplicação no Streamlit Cloud é realizado de forma contínua
(\textit{Continuous Deployment}): qualquer \textit{push} ao \textit{branch} principal
do repositório \texttt{caiorj2004/aih\_sp} no GitHub desencadeia automaticamente uma
nova compilação e publicação da aplicação. O ciclo de vida do ambiente de execução é
gerido integralmente pela plataforma, incluindo a instalação das dependências listadas
em \texttt{requirements.txt}.

A integração entre o GitHub e o Streamlit Cloud é estabelecida uma única vez mediante
autenticação OAuth, tornando o fluxo de publicação completamente transparente para o
desenvolvedor: a modificação de qualquer ficheiro do repositório reflecte-se na
aplicação publicada sem intervenção manual adicional.

% =============================================================
\chapter{Resultados e Conclusão}
% =============================================================

A metodologia aplicada revelou-se um sucesso técnico em todas as suas etapas.
A transição de paradigmas de bancos de dados — utilizando a complacência e velocidade
do \textit{SQLite} como \textit{Staging Area} preliminar para o posterior enquadramento
na governança estrita e alta concorrência de leitura do SGBD \textit{PostgreSQL} —
demonstrou profundo rigor arquitetural. Os dados transitaram de CSVs caóticos do
DATASUS para um \textit{Data Warehouse} modelado e higienizado, com tipagem forte e
valores monetários blindados em \texttt{Numeric(20,4)}.

Sobre esta base relacional robusta, foi construída e publicada uma camada de
apresentação completa na forma de um \textit{dashboard} web interativo. A escolha do
Streamlit como \textit{framework} permitiu condensar, num código de dimensão reduzida,
funcionalidades que noutras \textit{stacks} exigiriam a implementação de um
\textit{backend} REST separado, uma aplicação \textit{frontend} em JavaScript e uma
camada de autenticação independente.

A introdução do módulo \texttt{fallback.py} com dados Parquet locais revelou-se um
investimento de resiliência essencial: a aplicação continua operacional mesmo na
ausência de conectividade ao servidor institucional, degradando graciosamente apenas as
funcionalidades dependentes da dimensão UF. O toggle manual de modo de conexão confere
flexibilidade adicional para demonstrações e validações offline.

As otimizações implementadas no acesso ao banco de dados — em particular o padrão de
\textit{Loose Index Scan} via CTE recursiva, o uso de cursores \textit{server-side}
com \textit{streaming} em lotes e a decomposição da consulta consolidada em dois
estágios — demonstraram ser abordagens essenciais para garantir tempos de resposta
aceitáveis sobre tabelas de alta cardinalidade, sem necessidade de desnormalização ou
materialização de vistas pré-calculadas.

A secção de gráficos analíticos evoluiu de quatro para cinco visualizações,
incorporando o heatmap sazonal para identificação de padrões temporais e substituindo
o gráfico de donut por um treemap interativo com agrupamento de ``Outros''. O scatter
plot ganhou suporte ao coeficiente de Spearman para robustez a dados com distribuição
assimétrica — característica frequente em dados de procedimentos hospitalares, onde
poucos municípios concentram volumes muito elevados.

O conjunto do projeto — desde a extração automatizada no DATASUS até à publicação
do \textit{dashboard} analítico — constitui um pipeline de dados de ponta-a-ponta
completo, totalmente reprodutível e pronto para integração com ferramentas avançadas
de \textit{Business Intelligence}.

\end{document}
